\subsection{PERIODIC FUNCTIONS}

DEFINITION 1. A function \(f\) , defined on \(\mathbf{R}^n\) , with values in an arbitrary set \(\mathbf{E}\) , is said to be periodic if there exists a point \(a \neq 0\) in \(\mathbf{R}^n\) such that

\[
f (\boldsymbol {x} + \boldsymbol {a}) = f (\boldsymbol {x})\tag{1}
\]

for all \(x \in R^{n}\) . If f is periodic, every point \(a \in R^{n}\) for which the relation (1) is an identity in x is called a period of f.

The set G of all periods of a periodic function f is clearly a subgroup (which by hypothesis does not consist of o alone) of the additive group \(R^{n}\) . If f is a continuous periodic mapping of \(R^{n}\) into a Hausdorff topological space E, its group of periods G is closed. For if \(G_{x}\) denotes the set of all \(a \in R^{n}\) such that \(f(x + a) = f(x)\) for a given point \(x \in R^{n}\) , then G is the intersection of the \(G_{x}\) as x runs through \(R^{n}\) , and each \(G_{x}\) is closed (Chapter I, § 8, no. 1, Proposition 2). Let V be the largest vector subspace contained in G (no. 2, Theorem 2); the function f is constant on every coset mod V; if W is a vector subspace complementary to V, then f is determined by its restriction to W. In other words (W being a topological group isomorphic to \(R^{p}\) for some p), the study of continuous periodic functions on \(R^{n}\) is reduced to the study of such functions whose group of periods G is discrete; if this group is of rank q, the function f is said to be q-ply periodic, and every free system of q points which generate G is called a principal system of periods of f.

If \((\boldsymbol{a}_{i})\) and \((\boldsymbol{b}_{i})\) are two principal systems of periods of f, we have seen (no. 1) that each can be obtained from the other by a linear transformation with integer coefficients and determinant \(\pm1\) .

Let \(\varphi\) be the canonical mapping of \(\mathbf{R}^n\) onto \(\mathbf{R}^n/\mathbf{G}\) ; to every mapping \(g\) of \(\mathbf{R}^n/\mathbf{G}\) into a set \(\mathbf{E}\) corresponds the function \(\dot{g}=g\circ\varphi\) , which is a periodic mapping of \(\mathbf{R}^n\) into \(\mathbf{E}\) , having a group of periods which contains G; and conversely every mapping of \(R^{n}\) into E which has a group of periods containing G is of this form, since it is compatible with the relation \(x \equiv y \pmod{G}\) (Set Theory, R, § 5, no. 7). In this way we define a bijective mapping \(g \to \dot{g}\) of the set of all mappings of \(R^{n}/G\) into E onto the set of all mappings of \(R^{n}\) into E whose group of periods contains G. For \(\dot{g}\) to be continuous (E being a topological space) it is necessary and sufficient that g should be continuous (Chapter I, § 3, no. 4, Proposition 6).

